\originalchapter{复合系统、纠缠与全同粒子}
\originalsection{张量积和纠缠}
对两枚自旋, 联合测量有$++,+-,-+,--$四种结果; 每种配一个复振幅$c_{ij}$. 独立乘积制备时$c_{ij}=a_i b_j$, 一般联合制备却不必能这样分解. 整理这四个组合所需的联合态空间不是把维数相加, 而是张量积$\mathcal H_A\otimes\mathcal H_B$. 若两侧基为$\ket{i}_A$和$\ket{j}_B$, 联合基为$\ket{i}_A\otimes\ket{j}_B$, 一般归一态$\ket\Psi=\sum_{ij}c_{ij}\ket{i,j}$, $\sum_{ij}|c_{ij}|^2=1$. 两个二态系统有四个基向量, 四个复系数在归一化后描述联合纯态.

只作用在$A$上的算符为$A\otimes I_B$, 它在不同$B$分量中使用同一个$A$, 即$(A\otimes I)\ket{i,j}=(A\ket i)\otimes\ket j$. 例如$S_z\otimes I$在$\ket{+,-}$读$+\hbar/2$, 不改第二标签. 此处下标$A$也作系统名, 算符$A$是一个测量规则, 应由上下文分清. 若不存在$\ket\alpha\otimes\ket\beta$的乘积表达, 纯态称为纠缠态. 例如$\ket\Psi=(\ket{00}+\ket{11})/\sqrt2$. 假定它能写为$(a\ket0+b\ket1)\otimes(c\ket0+d\ket1)$, 则要求$ac,bd$非零而$ad,bc$为零, 彼此矛盾. 因而它不能分解.

在该态中测量$A$的$z$基, 结果0和1概率各半; 给定$A$结果0, $B$结果确定为0. 这说明关联, 并不允许$A$主动选择其随机测量结果来发送信息. 忽略$A$的结果时, $B$仍有相同的一半一半分布.

\originalsection{约化密度算符}
先区别两种同样在$z$上各半的制备: 随机选$z+$或$z-$给矩阵$\rho_{\rm mix}=\left(\begin{smallmatrix}1/2&0\\0&1/2\end{smallmatrix}\right)$; 相干$x+$给$\rho_{x+}=\left(\begin{smallmatrix}1/2&1/2\\1/2&1/2\end{smallmatrix}\right)$. 在$x$测量前者仍各半、后者必为$+$, 非对角项保留了相位干涉信息. 经典概率$p_s$不是可相加的复概率幅. 记$\Tr$为在任意正交归一基上的对角矩阵元之和, 称为迹. 纯态密度算符是$\rho=\ket\psi\bra\psi$, 统计混合则为$\rho=\sum_s p_s\ket{\psi_s}\bra{\psi_s}$, $p_s\ge0$, $\sum_sp_s=1$. $\rho\geq0$指每个矢的$\bra v\rho\ket v\geq0$, 不是要求每个矩阵元非负. 它满足$\rho\ge0$, $\Tr\rho=1$, 期望值写为$\Tr(\rho A)$. 混合态表示对制备结果或其他子系统未保留的信息, 一般不能用单个波函数代替.

不查看系统$B$的结果时, 把所有$B$标签的可能性求和, 仍留下$A$上的矩阵; 这称偏迹, 全迹则只给一个数. 将$\ket\Psi\bra\Psi=\sum_{iji\prime j\prime}c_{ij}c_{i\prime j\prime}^*\ket i\bra{i\prime}\otimes\ket j\bra{j\prime}$展开, 对$B$取迹$\Tr(\ket j\bra{j\prime})=\delta_{jj\prime}$, 就保留相同$B$标签的项. 联合纯态中只观测$A$时, 定义$\rho_A=\Tr_B\rho_{AB}$, 矩阵元为$(\rho_A)_{ii'}=\sum_j c_{ij}c_{i'j}^*$. 这样$\bra\Psi A\otimes I\ket\Psi=\Tr(\rho_AA)$. 对上面的纠缠态, $\rho_A=I/2$, 局部$S_z$均值0, $S_z^2$均值$\hbar^2/4$, 局部$x$也各半. 虽然联合态纯, 单个子系统的约化态混合. 相反,$(\ket0+\ket1)/\sqrt2$的纯态密度矩阵有非零非对角元, 与$I/2$不同; 在$x$基测量即可区别.

\originalsection{交换对称性}
电子是费米子, 光子是玻色子. 本节$x_i$的空间坐标与自旋标签在完整态中一同交换. 全同粒子的内禀属性相同, 用标签1、2不能增加可测的物理差别. 交换算符$P_{12}$把$\psi(x_1,x_2)$变为$\psi(x_2,x_1)$, 满足$P_{12}^2=I$. 标准三维非相对论量子理论中, 玻色子用对称态, 费米子用反对称态. 自旋与统计的联系是更深的相对论量子场论结论; 本册把相应交换规则作为物理输入, 不从$P^2=I$独自证明它.

两个正交单粒子态$a,b$的规范化联合态为
\[
\ket{a,b}_\pm=\frac{\ket a\otimes\ket b\pm\ket b\otimes\ket a}{\sqrt2}.
\]
若$a=b$, 反对称式为零, 无法归一化, 因而两个全同费米子不能占据同一个\emph{完整单粒子态}. 空间波函数相同但自旋不同, 不属于同一完整态, 可以占据同一空间轨道. 对电子的总态, 空间对称态须与反对称自旋单态搭配, 空间反对称态须与对称自旋三重态搭配, 使总态反对称.

\begin{example}
两个无相互作用、无自旋的全同粒子在同一无限深势阱中. 单粒子能量为$E_n=n^2E_1$. 比较两玻色子与两费米子的最低总能量.
\end{example}
\begin{solution}
玻色子可共同占据$n=1$, 所以总能量$2E_1$. 无自旋费米子必须占据两个不同单粒子态, 最低选择$n=1,2$, 反对称归一态为$(u_1(x_1)u_2(x_2)-u_2(x_1)u_1(x_2))/\sqrt2$, 总能量$5E_1$. “无自旋”是本数学模型的限制; 对真实自旋$1/2$电子, 可以用相反自旋占据同一空间基态, 最低能量就不同.
\end{solution}
